\subsection{Gravitational Shadows and the Spectral Wake}
  \label{sec:gravitational-shadows}

  Within the torsional framework described above,
  the constraint $\Omega_w$ induces an extended deformation of the
  local spectral structure of the $\chi$ substrate.
  This deformation modifies the local density of admissible relaxation
  modes and may persist over scales larger than the core configuration.

  We refer to this persistent redistribution of spectral capacity as a
  \emph{spectral wake} or gravitational shadow.
  It does not correspond to additional matter or propagating degrees of
  freedom, but to a long-lived modification of the relaxation structure
  induced by the torsional configuration.

  Such spectral wakes provide a possible structural interpretation for
  phenomena commonly attributed to dark matter, without requiring the
  introduction of new particle species.

  Two qualitative features naturally emerge.

  \paragraph*{Elastic remanence.}
    The spectral deformation may persist under displacement of the
    baryonic configuration.
    In this regime the wake behaves as a form of structural memory of the
    substrate, which could account for the observed offsets between
    baryonic and effective gravitational components in systems such as
    the Bullet Cluster.

  \paragraph*{Non-local susceptibility.}
    When the local relaxation gradient becomes sufficiently weak,
    the substrate response may transition from a linear to a nonlinear
    regime.
    A characteristic acceleration scale of order
    \(
    a_0 \sim c H_0
    \)
    then appears as the natural scale at which this transition occurs,
    leading to phenomenology similar to MOND-like modifications of
    gravitational dynamics.

    Within the Cosmochrony framework, such effects would therefore
    reflect a persistent redistribution of relaxation capacity in the
    relational substrate rather than the presence of unseen matter.
